\documentclass[11pt,a4paper]{article}
\newcommand{\fichetype}{Exercices}
% Thème A4 commun, librement modifiable.
\usepackage[T1]{fontenc}
\usepackage[utf8]{inputenc}
\usepackage[french]{babel}
\usepackage{lmodern,amsmath,amssymb,mathtools,tikz,pgfplots}
\usepackage[margin=2.2cm,headheight=15pt]{geometry}
\usepackage{xcolor,fancyhdr,enumitem,hyperref,needspace}
\usetikzlibrary{arrows.meta,positioning,calc}
\pgfplotsset{compat=1.18,width=11cm,height=6cm}
\definecolor{Ink}{HTML}{202124}
\definecolor{Blue}{HTML}{356B9E}
\definecolor{Muted}{HTML}{737579}
\color{Ink}
\renewcommand{\familydefault}{\sfdefault}
\setlength{\parindent}{0pt}
\setlength{\parskip}{6pt}
\setlist{itemsep=3pt,topsep=4pt}
\pagestyle{fancy}\fancyhf{}
\fancyhead[L]{\small\color{Muted}MATHS LYCÉE}
\fancyhead[R]{\small\color{Muted}2026--2027}
\fancyfoot[L]{\small\color{Muted}\fichetype}
\fancyfoot[R]{\small\thepage}
\renewcommand{\headrulewidth}{0.3pt}
\hypersetup{colorlinks=true,urlcolor=Blue,linkcolor=Blue}
\newcommand{\R}{\mathbb{R}}
\newcommand{\N}{\mathbb{N}}
\newcommand{\E}{\mathbb{E}}
\newcommand{\V}{\mathrm{V}}
\newcommand{\exercice}[3]{\Needspace{6\baselineskip}\vspace{0.5cm}{\color{Blue}\small\textbf{EXERCICE #1}\hfill #2}\par{\large\bfseries #3}\par}

\begin{document}
{\small\color{Blue}PREMIÈRE SPÉCIALITÉ}\par{\LARGE\bfseries Second degré : formes, équations et signes}\par{\large Exercices}\par\vspace{0.2cm}\hrule\vspace{0.4cm}
Les exercices progressent de l’application directe vers la recherche. Détailler les calculs et justifier les réponses.
\par\noindent\begin{minipage}{\linewidth}\exercice{1}{Application}{Reconnaître et convertir}
Pour $f(x)=x^2-8x+7$, déterminer les coefficients, compléter le carré et donner le sommet. En déduire le minimum.
\end{minipage}\par
\par\noindent\begin{minipage}{\linewidth}\exercice{2}{Application}{Trois discriminants}
Résoudre dans $\mathbb R$ : $x^2-7x+10=0$ ; $4x^2+4x+1=0$ ; $x^2+2x+5=0$.
\end{minipage}\par
\par\noindent\begin{minipage}{\linewidth}\exercice{3}{Entraînement}{Choisir la factorisation}
Factoriser $9x^2-16$, $x^2+x-12$ et $3x^2-12x$. Donner leurs racines.
\end{minipage}\par
\par\noindent\begin{minipage}{\linewidth}\exercice{4}{Entraînement}{Signe et inéquations}
Résoudre $2x^2-5x-3\le0$, puis $-x^2+2x+3<0$.
\end{minipage}\par
\par\noindent\begin{minipage}{\linewidth}\exercice{5}{Entraînement}{Une parabole imposée}
Un trinôme s’annule en $2$ et $5$ et vaut $-20$ en $0$. Déterminer son expression. Calculer son extremum.
\end{minipage}\par
\par\noindent\begin{minipage}{\linewidth}\exercice{6}{Entraînement}{Clôturer un jardin}
On dispose de $24$ m de grillage pour trois côtés d’un rectangle, le quatrième bordant un mur. Noter $x$ la longueur d’un côté perpendiculaire au mur. Exprimer l’aire et trouver les dimensions optimales.
\end{minipage}\par
\par\noindent\begin{minipage}{\linewidth}\exercice{7}{Entraînement}{Un algorithme de résolution}
Décrire un algorithme qui reçoit $a,b,c$ avec $a\ne0$ et renvoie le nombre de racines réelles et leurs valeurs. L’appliquer à $a=1,b=0,c=-2$. Pourquoi traiter $\Delta=0$ séparément ?
\end{minipage}\par
\par\noindent\begin{minipage}{\linewidth}\exercice{8}{Approfondissement}{Un paramètre}
Pour quels réels $m$ le trinôme $x^2-2mx+4$ est-il strictement positif pour tout réel $x$ ? Étudier les cas limites.
\end{minipage}\par
\end{document}
